From \parencite{park2026chain_of_retrieval}

Begrensning ved abstracts bruk:

Selve kritikken av abstract-basert gjenfinning:
"abstracts offer only sparse and highlevel summaries"
Det empiriske funnet (abstract-baselinjen slås av fulltekst-tilnærmingen):
"COR surpasses abstract-to-abstract (A2A) baselines by an average of 6.37% in Recall, demonstrating that simply relying on abstracts is suboptimal"
Et par merknader for hvordan du bruker dem.
Det første sitatet er deres generelle påstand om hva abstracts mangler. Det andre er deres måling. Til sammen dekker de både argumentet og beviset, så du trenger ikke flere.
Husk konteksten når du gjengir dette: funnet gjelder paper-til-paper-gjenfinning, der en hel artikkel er spørringen og fulltekst finnes på begge sider. Det er ikke direkte overførbart til din setting, der spørringen er en kort problembeskrivelse. Så bruk sitatene til å erkjenne en begrensning ved abstract-basert gjenfinning generelt og til å motivere fulltekst som videre arbeid, men ikke fremstill det som at de har vist at din tilnærming er underlegen, siden de ikke testet ditt scenario.

